%=============================================================================%
%  正文：福州大学开题报告五段式（一本 / 二研究内容与目标 / 三方案与可行性 /
%        四特色与创新 / 五参考文献）
%
%  写法约定
%    · 五个一级标题是学校表格里固定的，已用 \section{} 写好，不要增删、不要手打「一、」；
%      编号由模板自动生成（\chinese{section}、）。
%    · 以 % 开头的行是"填写要点"，编译时不显示；写完后留着当备忘、或者删掉都行。
%    · 需要三级标题时用 \subsubsection{}，图、表、公式、引用的写法见文末「写法示例」。
%=============================================================================%

\section{论文选题依据}

\subsection{选题背景与意义}
% 【写什么】3~4 段，只回答"为什么非做不可"，建议按这个顺序展开：
%   ① 领域/产业背景：这个测量或重建需求从哪来，有什么国家或行业需求（可先引 1~2 篇综述）；
%   ② 现有方法在本课题场景下失效的具体表现：要写现象和证据，不要只写"效果不佳"；
%   ③ 本课题的切入点，以及做成之后的价值（学术价值一句、工程价值一句）。
% 【注意】这一节不要写成文献综述，综述放到 1.2，否则两节会重复。

\subsection{国内外本学科领域的发展现状与趋势}
% 【写什么】按"技术方向"拆成 2~4 个三级标题，每个方向独立成段，段内顺序固定：
%   国外代表性方法（时间线 + 关键做法 + 关键指标） → 国内进展 → 这个方向还差什么。
%   最后回到本课题：上面这些"还差什么"，正是本课题要补的空白。
%
% 【要不要三级标题】要的话照下面这行改写，方向名自己起；不要的话保持注释即可。
% \subsubsection{方向一：×××的研究现状}
%
% 【引用】正文用 \cite{key} 引用，key 在 bib/ref.bib 里定义，上标编号由模板自动处理。


\section{论文的研究内容、研究目标及拟解决的关键问题}

\subsection{研究内容}
% 【写什么】分 3~4 条写，每条一句话说清"做什么"，不写方法细节（细节放第三章）。
%   常用写法：
%   \begin{enumerate}
%     \item 研究×××方法，解决×××问题；
%   \end{enumerate}

\subsection{研究目标}
% 【写什么】给可考核的目标：在什么数据集/场景上、达到什么精度或效率、相对现有基线提升多少。
%   反例："提高鲁棒性""实现高精度重建"——没法验收，评委会问"多少算高"。

\subsection{拟解决的关键问题}
% 【写什么】2~3 条，每条两句话：第一句说"难在哪"，第二句说"打算怎么破"。
%   这一节是评委判断你有没有想清楚的地方，宁可少写一条，也不要写空话。

% 【提示】若学院要求给出量化指标，可在 2.2 后面加一小节列张表：
%   指标名称 / 现有方法 / 本课题目标 / 验证方式。表格写法见文末「写法示例」。


\section{拟采取的研究方案及可行性分析}

\subsection{基本思路}
% 【写什么】一段话把技术路线串成因果链：输入是什么 → 中间表示是什么 → 输出是什么，
%   并说明"为什么这样串"（每一环为什么必须在前一环之后）。

\subsection{技术方案与技术路线}
% 【写什么】按研究内容拆成 3.2.1、3.2.2 … 每小节先"方案"（怎么做）再"路线"（步骤怎么走），
%   配 1~2 张流程图。方案里出现的关键公式要编号，后面正文要能引用它。

\subsection{可行性分析：研究条件与基础}
\subsubsection{已有研究的硬件基础}
% 【写什么】设备清单：型号 + 关键指标（分辨率、帧率、精度、量程），说明为什么够用。
\subsubsection{已有研究的实验基础}
% 【写什么】已经做完的事：标定结果、采集了多少组数据、跑通了什么流程、有什么初步结果。
%   这一节最好放真实数字和真实图，是评委会追问、也最容易加分的地方。

\subsection{研究进度安排}
% 【写什么】按四个阶段写，每阶段 2~4 条任务，动词开头（搭建/实现/验证/撰写）。
%   阶段之间要能接上，最后一段一定落到"写学位论文 + 答辩"。
本研究工作计划与进度时间节点安排如下所示：

\textbf{第一阶段（××××.××.××--××××.××.××）}
\begin{itemize}
	\item ×××
\end{itemize}

\textbf{第二阶段（××××.××.××--××××.××.××）}
\begin{itemize}
	\item ×××
\end{itemize}

\textbf{第三阶段（××××.××.××--××××.××.××）}
\begin{itemize}
	\item ×××
\end{itemize}

\textbf{第四阶段（××××.××.××--××××.××.××）}
\begin{itemize}
	\item ×××
\end{itemize}


\section{本课题的特色与创新之处}
% 【写什么】3~4 条，每条两段：先一句小标题，再一段"别人怎么做 → 本课题怎么做 → 差别在哪"。
%   创新点要能对应到前面的研究内容，不要凭空多出一条前面没做的事。
%
% 【写法（照结构，别照抄内容）】
% （1）×××
% 现有方法……；本课题……；因此……。

\proposalstep{（1）×××}
% 把上面的 \proposalstep 换成自己的小标题，逐条往下写。
% 若不想手动编号，也可以直接用 \paragraph{（1）×××} 或 \textbf{（1）×××}。

\proposalstep{（2）×××}

\proposalstep{（3）×××}


%=============================================================================%
%  写法示例（需要时把对应注释块复制到正文位置，再把内容换掉）
%
%  【插图】
%  \begin{figure}[H]
%    \centering
%    \includegraphics[width=0.95\textwidth]{images/proposal/你的图.jpg}
%    \caption{技术路线图}
%    \label{fig:route}
%  \end{figure}
%  正文引用：如图~\ref{fig:route}所示。      % 图号自动为「图 3-1」，按章编号
%
%  【表格】
%  \begin{table}[H]
%    \centering
%    \caption{实验条件对比}
%    \label{tab:cond}
%    \begin{tabular}{lccc}
%      \hline
%      项目 & 现有方法 & 本课题 & 备注 \\ \hline
%      ××× & ××× & ××× & ××× \\ \hline
%    \end{tabular}
%  \end{table}
%  正文引用：如表~\ref{tab:cond}所示。
%
%  【公式】
%  \begin{equation}
%    \label{eq:displacement}
%    \mathbf{d}(x) = \arg\min_{\mathbf{d}} \left\| \nabla I(x+\mathbf{d}) - \nabla I_0(x) \right\|_2^2
%  \end{equation}
%  正文引用：由式~\eqref{eq:displacement}可得……   % 公式连续编号，不带章号
%
%  【分条列举】itemize 用圆点，enumerate 用阿拉伯数字，两者都支持 \item
%
%  【加粗小标题】\textbf{×××} 或 \proposalstep{×××}（后者独立成段、不与下段分离）
%
%  【参考文献】条目写在 bib/ref.bib，正文 \cite{key} 引用；
%              bib/ref.bib 里现在只有注释掉的示例条目，取消注释、改成自己的即可。
%              注意：参考文献列表只显示正文引用过的条目，没引用就不会出现。
%=============================================================================%
